% Responsável: P1 (revisão final: todos)
\section{Introdução}\label{sec:intro}

\paragraph{O problema da organização.}
Uma equipa de verificação de factos recebe diariamente mais notícias do que consegue verificar.
Precisa de uma ferramenta que \emph{ordene} as notícias por suspeição, para que os verificadores
humanos comecem pelas mais prováveis de serem falsas. A ferramenta não decide o que é falso: apoia
a triagem. Esta distinção condiciona a escolha de métricas (o custo de um falso negativo e de um
falso positivo) e as considerações éticas (Secção~\ref{sec:etica}).

\paragraph{Pergunta de investigação.}
\emph{Que características linguísticas distinguem as notícias falsas das verdadeiras num corpus real
em português, e um detetor treinado com elas mantém o desempenho quando (a) se neutraliza o
comprimento dos textos e (b) é aplicado a temas que não viu no treino?}

\paragraph{Critério de sucesso (fixado antes de usar o teste).}
A pergunta considera-se respondida se: (i) identificarmos atributos com diferenças entre classes
estatisticamente significativas e interpretáveis; (ii) um modelo baseado no conteúdo superar a
\emph{baseline} que usa apenas o comprimento, com intervalos de confiança a 95\% não sobrepostos;
e (iii) quantificarmos a queda de desempenho em temas não vistos face ao desempenho no mesmo tema.

\paragraph{Contribuições.}
\begin{enumerate}
\item Um protocolo de avaliação para o Fake.Br sem fugas de informação: divisão agrupada por par de
notícias e por textos repetidos entre pares (Secção~\ref{sec:config}).
\item A evidência de que o comprimento, sozinho, explica a maior parte do desempenho reportado em
textos completos, e uma comparação de representações quando esse atalho é removido.
\item Uma análise convergente com três tarefas (classificação, tópicos e extração de entidades) que
aponta o \emph{estilo} e a \emph{ancoragem factual} (datas, locais) como os sinais robustos, e não o tema.
\item Um repositório que reproduz a tabela principal com um único comando: \repo.
\end{enumerate}

O resto do artigo está organizado assim: a Secção~\ref{sec:rel} discute o trabalho relacionado; a
\ref{sec:dados} caracteriza o corpus; as Secções~\ref{sec:metodo} e~\ref{sec:config} descrevem o
método e a configuração experimental; a \ref{sec:resultados} apresenta os resultados, discutidos na
\ref{sec:discussao}; a \ref{sec:etica} trata as questões éticas e as limitações, e a
\ref{sec:conclusao} conclui.
